\documentclass[11pt,letterpaper, twocolumn]{article}


\setlength\parindent{0pt}% indent width
\usepackage[utf8]{inputenc} %paquete que permite colocar acentos
\usepackage{helvet} %The package pro­vides a means of loading Helvetica (and its clones), scaled so that it matches Times Ro­man (to some ex­tent).
\renewcommand{\familydefault}{\sfdefault} % Fuente en Arial
\usepackage{hyperref} %Paquete para referenciar
\hypersetup{
    colorlinks=true,
    linkcolor=blue,
    filecolor=magenta,      
    urlcolor=blue,
}
\usepackage[spanish]{babel} %Paquete para el idioma español
\usepackage{amsmath}
\usepackage{amsfonts}
\usepackage{amssymb}
\usepackage{graphicx}
\usepackage{fancyhdr}%encabezados
\usepackage[left=2.5cm,right=2.5cm,top=2.5cm,bottom=2.5cm]{geometry}%tamaño márgenes
\usepackage{lmodern}
\providecommand{\abs}[1]{\lvert#1\rvert} %valor absoluto
\providecommand{\norm}[1]{\lVert#1\rVert} %norma
\renewcommand{\baselinestretch}{1} %interlineado
\usepackage[small]{caption}  %Tamaño del caption de las figuras

%---- encabezado -----
\pagestyle{fancy}
\fancyhf{}
\rhead{\thepage}
\lhead{\textcolor[cmyk]{1,0,1,0}{\underline{Laboratorio de Electromagnetismo}}}
\renewcommand{\headrulewidth}{0pt}
\usepackage{titling}


%-------Título----------------
\title{\protect\vspace{-3cm}\protect\includegraphics[width = 3cm]{Logos/logofisica.png}\protect\includegraphics[width = 6.5cm]{Logos/logouis.png}\\ ``TÍTULO DE LA PRÁCTICA DE LABORATORIO"  }
\author{Nombres y Apellidos del Autor 1. \and Nombres y Apellidos del Autor 2. \and Nombres y Apellidos del Autor 3.\\
\small{Universidad Industrial de Santander - Escuela de física}}
\date{} %ingresar la fecha o simplemente borrar los "{}"

\begin{document}

\twocolumn[
  \begin{@twocolumnfalse}
    \maketitle
    
%------------------------ Resumen----------------------------------
\begin{abstract} 

En el resumen se debe describir brevemente la experiencia y la principal conclusión de la práctica realizada, acorde con los objetivos de la práctica de laboratorio, con máximo de 300 palabras.  En el resumen se plantea el problema y los principales resultados obtenidos.
      
\end{abstract}
%----------------------------------------------------------------------------------------------------------------------------------------    
  \end{@twocolumnfalse}
  \vspace{0.7cm}
  ]
  
%----------------------------------------------------------------------------------------------------------------------------------------

\section{Introducción}
En la introducción se realiza planteamiento del problema, es decir, definir el objeto del informe: ¿por qué se realiza el estudio?, se hace una breve descripción de los conceptos pertinentes estado del arte para el tema de investigación con sus respectivas respectivas referencias. 

\section{objetivos}

Aquí se plantea el objetivo general de la práctica de laboratorio, junto con los objetivos específicos, si los hay, que fueron necesarios para completar o satisfacer el objetivo general

%----------------------------------------------------------------------------------------------------------------------------------------

\section{Marco teórico}

El marco teórico consiste en la recopilación de antecedentes, investigaciones previas y consideraciones teóricas por donde se sustenta un proyecto de investigación, análisis, hipótesis o experimento, permitiendo la interpretación de los resultados y la formulación de conclusiones. Es importante que se coloquen las referencias de los libros o guías de consulta, así como de las imágenes utilizadas (en caso de no ser hechas por usted)

%----------------------------------------------------------------------------------------------------------------------------------------

\section{Metodología}
La metodología responde a la pregunta: ¿cómo se hizo la investigación?  En esta parte del reporte de investigación se presenta una descripción detallada del procedimiento específico realizado en el desarrollo del proyecto de investigación mencionando los equipos utilizados. Se debe incluir un gráfico, dibujo o foto del montaje realizado en el desarrollo del experimento, titulado como se muestra al final de este documento.

%----------------------------------------------------------------------------------------------------------------------------------------

\section{Tratamiento de datos}
En esta sección se incluye la tabla de datos, los cálculos pertinentes y una tabla donde se resuman los resultados con los cálculos realizados, además del cálculo del error, en caso que sea necesario.

%----------------------------------------------------------------------------------------------------------------------------------------

\section{Análisis de resultados}

En esta sección se explican los resultados obtenidos y la comparación de dichos datos con los obtenidos por otros investigadores o los datos esperados (estos se deben referenciar, usualmente deben estar en el marco teórico), además, en lo posible, se deben incluir gráficos de análisis que muestren la relación entre las variables que hacen parte de la investigación.  Se recomienda primero comenzar con la identificación de las relaciones que los resultados indican.  Posteriormente, señalar los casos en los cuales no se encuentre correlación entre las variables.  De igual manera, se deben establecer los aspectos que no se pudieron resolver con el desarrollo de la experiencia, explicando las razones por las cuales se presentaron.  En este aspecto es importante no ocultar, alterar o forzar los datos, sino reconocer las limitaciones de los datos, de los instrumentos de medición y los posibles errores en el procedimiento de medición.

Finalmente, se recomienda exponer las consecuencias teóricas de la investigación y las posibles aplicaciones de la misma.  El análisis de resultado guía el planteamiento de las conclusiones, pero no repite textualmente.

%ejemplo para insertar imágenes o gráficos en latex

\begin{figure}[h]
    \centering
    \includegraphics[width = 2in]{Logos/logofisica.png}
    \caption{Ejemplo para colocar gráficas en \LaTeX. \textit{Imagen tomada de: \url{http://eno2017.ciencias.uis.edu.co/}}   }
    \label{ejemplo1}
\end{figure}

%----------------------------------------------------------------------------------------------------------------------------------------


\section{Conclusiones}

Aquí se describen las conclusiones de la investigación.  Las conclusiones deben tener mínimo 200 palabras y estar en coherencia con los objetivos de la investigación y  el análisis de resultados. En esta sección es necesario que el estudiante realice un balance sobre la comprensión lograda acerca de los aspectos teóricos y experimentales.

%---------------------------------------------------------------------------------------------------------------------------------------

\section{Referencias}

se deben referenciar las fuentes consultadas para el marco teórico, se recomienda en lo posible que el autor cree las gráficas en un programa de acceso libre como \href{https://inkscape.org/es/}{Inkscape} pero en el caso de tomar gráficas de otra fuente, esta se debe referenciar como en el caso de la figura \ref{ejemplo1}.

Las referencias se pueden hacer a mano, no obstante, existe una herramienta muy útil de latex para colocar referencias llamada \href{https://www.latex-tutorial.com/tutorials/bibtex/}{Bibtex}, se recomienda acceder al enlace y consultar el tutorial sobre cómo usar esta herramienta

%---------------------------------------------------------------------------------------------------------------------------------------

\newpage

\twocolumn[
\begin{@twocolumnfalse}
\section{Anexos}
En esta sección deben ir las fotografías o imágenes escaneadas de las tablas de datos que registró en hoja de trabajo que usó durante la sesión práctica en el laboratorio, que incluye la fecha y firma de su profesor de Laboratorio.
\end{@twocolumnfalse}
]

%----------------------------------------------------------------------------------------------------------------------------------------
\end{document}
